\documentclass[showpacs,amsmath,amssymb,twocolumn,superscriptaddress,notitlepage,preprintnumbers,pra]{revtex4-1}
% \usepackage{ctex}
\usepackage[dvips]{graphicx}
\usepackage{amsmath,amssymb,amsthm,mathrsfs,amsfonts,dsfont}
\usepackage{subfigure, epsfig}
\usepackage{braket}
\usepackage{bm}
\usepackage{enumerate}
\usepackage{color,xcolor}
\usepackage{qcircuit}
\usepackage{graphicx}
\usepackage{algorithm}
\usepackage{algpseudocode}
\usepackage{algpseudocode}
\usepackage{hyperref}
\hypersetup{colorlinks=true, linkcolor=blue, citecolor=blue, urlcolor=black }
\usepackage{comment}
\usepackage[normalem]{ulem}
 

\newcommand{\tr}{\operatorname{Tr}}
\newcommand{\Var}{\mathrm{Var}}
 
\newcommand{\Cohere}{\Gamma}
\DeclareMathOperator*{\argmax}{arg\,max}
\DeclareMathOperator*{\argmin}{arg\,min}

% \newcommand{\mc}{\mathcal}
\newcommand{\mb}{\mathbf}
\newcommand{\mbb}{\mathbb}
\newcommand{\mr}{\mathrm}
\newcommand{\abs}[1]{\lvert#1\rvert}


\newcommand{\sinc}{\mathrm{sinc}}
\newcommand{\sech}{\mathrm{sech}}

\newcommand{\Ht}{H$_2$ }
\newcommand{\mc}[1]{\mathcal{#1}}
%\newcommand{\hat}[1]{\mathcal{#1}}
\DeclareMathOperator*{\E}{\mathbb{E}}

\usepackage{varioref}
\labelformat{equation}{(#1)}
\def\equationautorefname{Eq.}
\labelformat{section}{#1}
\def\sectionautorefname{Sec.}
\labelformat{figure}{#1}
\def\figureautorefname{Fig.}
\labelformat{proposition}{#1}
\def\propositionautorefname{Proposition}
\labelformat{lemma}{#1}
\def\lemmaautorefname{Lemma}
\labelformat{theorem}{#1}
\def\theoremautorefname{Theorem}
\labelformat{observation}{#1}
\def\observationautorefname{Observation}
\labelformat{definition}{#1}
\def\definitionautorefname{Definition}
\labelformat{problem}{#1}
\def\problemautorefname{Problem}
\labelformat{algorithm}{#1}
\def\algorithmautorefname{Algorithm}
\labelformat{corollary}{#1}
\def\corollaryautorefname{Corollary}
\labelformat{statement}{#1}
\def\statementautorefname{Statement}

\newtheorem{theorem}{Theorem}
\newtheorem{observation}{Observation}
\newtheorem{proposition}{Proposition}
\newtheorem{statement}{Statement}
\newtheorem{lemma}{Lemma}
\newtheorem{remark}{Remark}
\newtheorem{condition}{Condition}
\newtheorem{corollary}{Corollary}
\newtheorem{definition}{Definition}
\newtheorem{assumption}{Assumption}
\newtheorem{problem}{Problem}

\newcommand{\sun}[1]{\textcolor{blue}{#1}}
\newcommand{\blue}{\textcolor{blue}} %blue text
 

\begin{document}
 
 

\title{Energy consumption estimates}
 

\date{\today}

 
% --------------------  ABSTRACT  --------------------

% ------------  AUTHORS AND AFFILIATIONS ----------
 

 

\begin{abstract}
 
 

\end{abstract}

\maketitle


 
\section{Introduction}
\subsection{(Proposed) Hardware Parameters:}

I make the following modelling assumptions for the IBM processors:

We use the following hardware parameters in the temperature dependent model:
\begin{itemize}
    \item $\gamma_0$ is fit at runtime to the known values of $T_1$ at the default operational mixing chamber temperature
    \begin{itemize}
        \item Assume the operational cryostat temperature to be 0.013K
    \end{itemize}
    \item $\gamma_{\varphi-\text{base}}$ is fit at runtime to the known values of $T_2$ at the default operational mixing chamber temperature
    \begin{itemize}
        \item Again, assume the operational cryostat temperature to be 0.013K 
    \end{itemize}
    \item $\frac{\gamma_{MXC}}{\gamma_0}$ is $0.85$.
    \begin{itemize}
        \item This represents the maximum value from the experimental dataset reported in Ref.~\cite{lvov2025thermometry}.
    \end{itemize}
    \item $T_{env}$ as 30mK
    \begin{itemize}
        \item Effective temperatures as low as 20mK have been observed Ref.~\cite{simbierowicz2024thermal_noise}.
    \end{itemize}
    \item $\Delta$ (superconducting energy gap) of $2.88e^{-23}$
    \begin{itemize}
        \item This is the known value for Aluminium, which we can assume the electrodes used in IBM transmon qubits are made of
    \end{itemize}
    \item $\omega_{ge}$ is extracted from known configuration data for each backend individually for each qubit
    \item $N_e$, the number of effective channels in the Josephson Junction, is estimated at $1.29\times10^{5}$. We calculate this value using:
    \begin{itemize}
        \item An estimated normal resistance of the Josephson Junction of $5\text{k}\Omega$
        \item Transparency of the junction in the subgap regime of $5\times10^{-5}$
    \end{itemize}
    \item $E_c$, charging energy of the capacitors in the transmon qubits, of $0.95\mu$eV
    \item $E_j$, josephson energy of the transmon qubits, of $62.1\mu$eV
\end{itemize}

Starting with an estimated value of $\gamma_0$: $5e3$, and an estimated standard operational cryostat temperature of 0.013K, we minimise the squared error
between the calculated and known, default relaxation time of the qubits. We then do the same for the value of $\gamma_{\varphi-\text{base}}$ keeping
$\gamma_0$ constant

The final scale factors for the value of $T_1$ and $T_2$ for each qubit are stored and applied to the calculated values to maintain granularity

When calculating gate errors, the difference between the known gate error and calculated gate error at the operational temperature is stored and assumed to be
representative of the non-thermal error, and thus independent of temperature. It is then added to future calculated gate errors

\bibliography{references.bib}
\end{document}	